%! TEX root = ../main.tex

\clearpage

\section{Thiết kế theo chức năng}

\subsection{Chức năng XX (XX: tên cụ thể) }
<Giải thích cách thức miền thông tin của hệ thống được chuyển sang các cấu trúc dữ liệu. Mô tả 
cách thức dữ liệu chính hay các thực thể của hệ thống được lưu trữ, được xử lý và được tổ chức. 
Liệt kê các cơ sở dữ liệu hay các mục lưu trữ dữ liệu.> 

\begin{itemize}
  \item \textbf{Mục đích:}
  \item \textbf{Giao diện:} hiển thị các ảnh giao diện từ góc nhìn của người sử dụng. Chúng có thể 
được vẽ bằng tay hay dùng công cụ vẽ tự động. Ta nên tạo ra chúng chính xác như có 
thể. Ta cũng có thể đánh số cho từng thành phần trong giao diện.  
  \item \textbf{Các thành phần trong giao diện:} ghi vào bảng sau các mô tả về từng thành phần (đã 
được đánh số) của giao diện. 

  \begin{tabularx}{\textwidth}{|R{1cm}|L{3.5cm}|L{4cm}|X|}
      \hline
      \rowcolor{gray!25} 
      \multicolumn{1}{|c|}{\textbf{STT}} & 
      \textbf{Loại điều kiện} & 
      \textbf{Giá trị mặc định} & 
      \textbf{Lưu ý} \\ \hline
      
      1 & 
      Mỗi thành phần trong giao diện có thể là button hay textbox hay combobox, v.v. & 
      a & 
      Viết lưu ý cho những thành phần trong giao diện có cách xử lý đặc biệt hoặc các quy định mà lập trình viên phải thực hiện. \\ \hline
  \end{tabularx}

  \item \textbf{Dữ liệu được sử dụng:}  liệt kê các bảng trong cơ sở dữ liệu hoặc các cấu trúc dữ liệu 
được cần đến bởi chức năng này. 

  \begin{tabularx}{\textwidth}{|R{1cm}|X|C{1.8cm}|C{1.8cm}|C{1.8cm}|C{2.2cm}|}
      \hline
      \rowcolor{gray!25} 
      % Dòng 1: Để trống 2 ô đầu, chỉ tạo header cho "Phương thức"
      & & \multicolumn{4}{c|}{\textbf{Phương thức}} \\ \cline{3-6}
      
      \rowcolor{gray!25}
      % Dòng 2: Khai báo multirow{-2} để đẩy chữ lên giữa 2 dòng và đè lên màu xám
      \multicolumn{1}{|c|}{\multirow{-2}{*}{\textbf{STT}}} & 
      \multicolumn{1}{c|}{\multirow{-2}{*}{\textbf{\begin{tabular}[c]{@{}c@{}}Tên bảng / \\ Cấu trúc dữ liệu\end{tabular}}}} & 
      \textbf{Thêm} & 
      \textbf{Sửa} & 
      \textbf{Xóa} & 
      \textbf{Truy vấn} \\ \hline
      
      % Dòng dữ liệu
      1 & Danh mục sách & x & x & x & x \\ \hline
      2 & Tài khoản người dùng & x & x & & x \\ \hline
  \end{tabularx}

  \item \textbf{Cách xử lý:} giải thích bằng lời hoặc vẽ sơ đồ mô tả dòng xử lý trên giao diện. 
  \item \textbf{Hàm/ sự kiện} (nếu có): mô tả giải thuật cho từng biến cố bằng sơ đồ hoặc bằng ngôn 
ngữ giả. 
  \item \textbf{Các ràng buộc (nếu có)}: ví dụ như tham khảo đặc tả nào của tài liệu đặc tả nào. 
\end{itemize}


\subsection{Chức năng YY (YY: tên cụ thể) }
<Có các mục tương tự như chức năng XX> 

